\scanpage{223}
\mspage{173}
% L1–5
\uline{NB} L'équivalence de a) et a$'$) devrait être triviale, mais pour cela il faut avoir défini les $\pi_n(X,x)$, pour $x$ variable dans $\mathrm{Ob}\,X$, comme système local sur $X$, chose que je n'ai pas fait encore. Je vais essayer de le faire dans un moment d'aise. 
% L5 margin
% slanted
\margin{\uline{NB} $W_{\unsure{\omega|\infty}}$ est le plus petit localiseur de $\mathrm{Cat}$ (satisfaisant \unsure{W$\mathscr L$|WL} 2$^\circ$ bis, b)).}
% L6–8
Tout d'abord, les $\Omega(X,x)$ sont les fibres d'une $W_\omega$-fibration parfaite, s'insérant dans le carré cartésien
% L9–11
\diagram{\begin{tikzcd}[ampersand replacement=\&,column sep=large]
\underline{\mathrm{Ch}}(X) \arrow[d,"{p_X=(\alpha_X,\beta_X)}"] \& \Omega(X,-) \arrow[l] \arrow[d]\\
X\times X \& X \arrow[l,"{\mathrm{diag}_X}"] \arrow[ul,dotted,"{i_X}"']
\end{tikzcd},}
% L12
On sait dès lors que l'on a un foncteur
% L13–14
\[
\begin{aligned}
\Pi X&\longrightarrow\mathrm{Hot}_{W_\omega}\\
x&\unsure{\mapsto|\vdash}\mathrm{hot}_{W_\omega}(\Omega(X,x))
\end{aligned}
\tag*{($\unsure{*|\alpha}$)}
\]
% L15
Or on a, pour $i\geq1$
% L16
\[
\pi_i(X,x)=\pi_{i-1}(\underbrace{\Omega(X,x)}_{\text{\unsure{objet} ponctué}}).
\]
% L17–20
Cela donne le système local cherché des $\pi_i$, si \unsure{seulement} on avait un foncteur un peu plus précis que ça, savoir
% L17–20 margin
% slanted
\margin{Mais il faut avoir vérifié que les $\pi_n(X,x)$ ($n\geq1$) sont des foncteurs \unsure{de|sur} $\mathrm{Hot}_{\bullet W_\infty}$, i.e. transforment $W_\infty$-équivalences en isom. Cela résulte \unsure{des} W(4\unsure{--|bis}) pour $W_\infty$, qui implique que $(X,x)\unsure{\longmapsto|\longrightarrow}\Omega(X,-)$ \unsure{est un foncteur} $\mathrm{Hot}_{\bullet W_\infty}$ \unsure{en lui-même|en fait} (\unsure{ceci est plus précisément} \illegible{2} \unsure{avec} $\mathrm{Hot}_{\bullet W_\infty}$).}
% L20
\[
\Pi X\longrightarrow\mathrm{Hot}_{\bullet\,W_\omega}
\qquad\text{(Catégorie homotopique \uline{ponctuée}).}
\]
% L21–22
Or la \deleted{application} flèche $i_X:X\unsure{\longrightarrow|-}\underline{\mathrm{Ch}}\,X$ \unsure{relative à|relèv. à} $p_X$ donne une section \inserted{can.} de $\Omega(X,-)$ sur $X$ \textendash
% archive sheet number: 223
